\section{Evaluation Design}\label{sec:b-design}

The evaluation asks four questions about the approach as a developer would use it.
\begin{itemize}
\item \textbf{RQ1, faithfulness.} Does the approach compute, without running the program, the circuit the program
builds, the slices of that circuit, and the classical statements the slices need?
\item \textbf{RQ2, coverage.} How many real programs can it slice, and why does it fail on the others?
\item \textbf{RQ3, cost.} How fast is it, and how does it scale against tools that work on the circuit?
\item \textbf{RQ4, what a slice shows.} How much of a program does a slice ask a developer to read, and what does it add
to the warnings of a bug detector?
\end{itemize}
Every analysis was planned in a protocol and changed only through dated deviations, which the replication package
contains, and all of them are exploratory.

\subsection{Subjects}
Table~\ref{tab:b-subjects} lists the programs and circuits. Bugs4Q~\cite{zhao2023bugs4q} and the corpus that LintQ's
authors collected from GitHub~\cite{paltenghi2024lintq} are real programs, most written for Qiskit releases before~1.0.
Qiskit HumanEval~\cite{vishwakarma2024qhe} holds tasks solved by Qiskit experts for current releases, and the Qiskit
textbook~\cite{qiskittextbook} programs teach the standard algorithms. We wrote three sets ourselves, the
\numToolDevPrograms{} runnable programs used to check the trace, programs in which a size constant grows the circuit
while the source stays fixed, and programs with loops that run on the device. The OpenQASM~2 front end is checked on
the \numSmlChecked{} OpenQASM~2 files up to 5~MB of MQT Bench~\cite{quetschlich2023mqt}, QASMBench and
VeriQBench~\cite{chen2022veriqbench}, and on test programs for the grammar and the specification's examples, and timed on the medium and large circuits of QASMBench~\cite{li2023qasmbench}.

\begin{table*}[t]
  \centering
  \caption{Programs and circuits of the evaluation}
  \label{tab:b-subjects}
  \footnotesize\setlength{\tabcolsep}{4pt}\renewcommand{\arraystretch}{1.15}
  \newcommand{\sep}{\arrayrulecolor{black!25}\hline\arrayrulecolor{black}}
  \begin{tabularx}{\textwidth}{@{}>{\raggedright\arraybackslash}X r r r r r >{\raggedright\arraybackslash}p{2.6cm}@{}}
    \toprule
    \textbf{Subject} & \textbf{Count} & \textbf{Traced} & \textbf{Qubits} & \textbf{Operations} & \textbf{Code lines} &
      \textbf{Used for} \\
    \midrule
    LintQ's corpus, GitHub programs & \numSubjLintqCount{} & \numSubjLintqTraced{} & \numSubjLintqMinQubits{}--\numSubjLintqMaxQubits{} &
      \numSubjLintqMedianOps{} / \numSubjLintqMaxOps{} & \numSubjLintqMedianLines{} & RQ1 to RQ4 \\ \sep
    Bugs4Q programs & \numSubjBugsCount{} & \numSubjBugsTraced{} & \numSubjBugsMinQubits{}--\numSubjBugsMaxQubits{} &
      \numSubjBugsMedianOps{} / \numSubjBugsMaxOps{} & \numSubjBugsMedianLines{} & RQ1, RQ2 \\ \sep
    Qiskit textbook & \numSubjTextCount{} & \numSubjTextTraced{} & \numSubjTextMinQubits{}--\numSubjTextMaxQubits{} &
      \numSubjTextMedianOps{} / \numSubjTextMaxOps{} & \numSubjTextMedianLines{} & RQ2 \\ \sep
    Qiskit HumanEval & \numSubjQheCount{} & \numSubjQheTraced{} & \numSubjQheMinQubits{}--\numSubjQheMaxQubits{} &
      \numSubjQheMedianOps{} / \numSubjQheMaxOps{} & \numSubjQheMedianLines{} & RQ2 \\ \sep
    Parametric programs, 5 sizes each, ours & \numSubjParamCount{} & \numSubjParamTraced{} of 20 &
      \numSubjParamMinQubits{}--\numSubjParamMaxQubits{} & \numSubjParamMedianOps{} / \numSubjParamMaxOps{} &
      \numSubjParamMedianLines{} & RQ3 \\ \sep
    Programs with device loops, ours & \numSubjCfCount{} & \numSubjCfTraced{} & \numSubjCfMinQubits{}--\numSubjCfMaxQubits{} &
      \numSubjCfMedianOps{} / \numSubjCfMaxOps{} & \numSubjCfMedianLines{} & RQ1 \\ \sep
    QASMBench medium and large, OpenQASM~2 & \numSubjQbCount{} & \numSubjQbTraced{} & \numSubjQbMinQubits{}--\numSubjQbMaxQubits{} &
      \numSubjQbMedianOps{} / \numSubjQbMaxOps{} & -- & RQ3 \\
    \bottomrule
  \end{tabularx}\\[3pt]
  \begin{minipage}{\textwidth}\footnotesize
  Traced counts the programs from which \tool{} reconstructs a circuit, the sizes of the parametric programs, and for
  QASMBench the circuits in valid OpenQASM~2. Qubits and operations describe the traced circuits, as smallest to largest
  and as median / largest, and \numToolMaxOps{} operations is the trace limit. Code lines exclude blank and comment
  lines and are medians. Programs by others are read and never run.
  \end{minipage}
\end{table*}

\subsection{References for faithfulness}
For RQ1, our runnable programs are also run with Qiskit, and the trace must equal the circuit Qiskit builds, operation
by operation. The slices must equal the backward light-cone slices of that built circuit, which the study computes with
code separate from the tool, and on the Bugs4Q bugs that port to circuits, the line slices must equal those obtained by
recording which line appends each operation while the port runs. The OpenQASM~2 front end's trace must equal the circuit
Qiskit loads from the same file, on every OpenQASM~2 file of the corpora, it must reject exactly the files Qiskit
rejects, and its srcML must give back the file byte for byte. For the programs with device loops, an exact density-matrix simulation that branches
over measurement outcomes deletes every operation outside a qubit's slice from the source and checks that the qubit's
state is unchanged at \numCfInputs{} inputs, and a negative control reads each device loop's body once. The classical
statements of E4 are compared with the values E2 tracks along the one path it follows, which give a second,
flow-sensitive set of lines, on our programs and Bugs4Q, on a sample of LintQ's corpus, and on a held-out sample of
\numSsHeldFiles{} files of that corpus that were not inspected while the srcSlice fixes were made. We fixed a
threshold of 95\% containment before the comparison.

\subsection{Coverage, cost and what a slice shows}
For RQ2, \tool{} traces every program, and the programs that yield no circuit are classified by a fixed rule on their
text. For RQ3, \tool{} slices every qubit, up to 64, of each program of LintQ's corpus that builds a circuit, with a limit
of 120~s per program. On the QASMBench circuits it slices every qubit in one pass, against three tools that slice one
qubit of the built circuit, the ancestors in Qiskit's DAG, Qiskit's \texttt{LightCone} pass and a reimplementation of the
QPDG~\cite{yu2025qpdg}, with a limit of 600~s per circuit for the baselines and 3{,}600~s for \tool{}, whose
srcML markup of the largest circuits takes minutes. On the parametric programs, a QFT, a GHZ preparation, a QAOA
ring and a repetition code, a size constant grows the trace while the source stays the same. All times are wall-clock on
one machine, an Intel Xeon W-2140B at 3.2~GHz with 32~GB of memory, with Python~3.12 and Qiskit~2.3. For RQ4, we measure
the share of a program's code lines that a qubit's slice holds, and compare \tool{} with LintQ on LintQ's public data,
without running LintQ. For each warning of LintQ's five analyses that concern qubits and bits, \tool{} traces the
program and evaluates a predicate on its qubit profiles, fixed before the comparison, at the warned line, on the warnings
LintQ's authors labelled true and false positives and on all warnings of these analyses. The examples in the papers of LintQ
and QChecker~\cite{paltenghi2024lintq, zhao2023qchecker} are transcribed as programs and sliced.
